% This is a template for your written document.
%
% To compile using latexmk on the command line, run the following: 
% latexmk -pdf main.tex

\documentclass[12pt]{article}
\usepackage{setspace}
\singlespace
\usepackage[left=1in,right=1in,top=1in,bottom=1in]{geometry}
\usepackage{graphicx}

\title{\textbf{A Software-Based Comparative Study of Error Sensitivity in Classical and Quantum Optimization Algorithms}}
\author{Aayush Rajak}

\begin{document}

\maketitle

\section{Introduction}
%The introduction will broadly introduce your topic, motivate it, and describe what will be done.
Everyone is familiar with the idea of a classical computer, but not everyone is familiar with the idea of a quantum computer. 
The Classical computing processes information using bits that take definite values of 0 or 1, whereas quantum computing 
uses qubits that can exist in superpositions of states and become entangled, enabling fundamentally different computational behavior
~\cite{wong2019}. Quantum computing has the potential to improve the performance of certain computational tasks, particularly search 
and optimization problems. Many quantum algorithms promise theoretical speedups over classical methods; however, these results often rely on 
ideal assumptions that ignore noise, limited circuit depth, and other practical constraints of current quantum technology. As a result, 
there is a gap between theoretical expectations and the real-world behavior of quantum algorithms. This work aims to bridge that gap through 
a software-based comparison of classical and quantum optimization algorithms, focusing on their sensitivity to noise and errors using 
classical simulation.


This research evaluates quantum optimization algorithms under realistic execution conditions and compares them with classical heuristics. 
Rather than focusing on theoretical speedups, it examines practical performance, robustness, and solution quality in noisy environments. 
Classical methods—genetic algorithms and greedy search—are compared with quantum approaches including Grover-based 
search and QAOA. A modular benchmarking framework with configurable noise models enables fair, reproducible 
evaluation using metrics such as solution quality, stability, circuit depth, and noise sensitivity. Results show quantum performance often 
degrades under realistic constraints, while classical heuristics remain more robust, highlighting the importance of hybrid strategies and 
practical evaluation methods.


This paper is organized into five benchmark comparisons used to evaluate classical and quantum optimization algorithms across representative 
problem types. Each row pairs a classical algorithm with a quantum counterpart and identifies the problem instance used to compare them. The 
first and primary comparison evaluates a Genetic Algorithm against QAOA on the Max-Cut problem, a standard benchmark for quantum optimization 
that allows a fair global search comparison. The second compares a Greedy Algorithm with QAOA on the Knapsack problem, highlighting the 
difference between a simple heuristic and a variational quantum method for constrained optimization. The third experiment contrasts a Genetic 
Algorithm with Grover’s Algorithm using MAX-SAT formulated as a search problem, emphasizing heuristic versus quadratic-speedup search. 
The fourth compares Greedy search with Grover’s Algorithm on Number Partitioning, illustrating structured versus unstructured search. 
Finally, a combined Genetic+Greedy versus QAOA+Grover evaluation on Random QUBO provides a unified benchmark across all algorithms.

The core hypothesis of this work is that while quantum algorithms offer theoretical advantages, their practical performance in near-term 
optimization tasks is inversely proportional to the level of environmental noise and circuit complexity. Specifically, it is hypothesized 
that classical heuristics—such as Genetic Algorithms and Greedy Search—will demonstrate superior robustness and solution quality compared to 
their quantum counterparts (QAOA and Grover’s Algorithm) when executed under realistic, noisy simulation parameters.

Furthermore, the research posits that a "cross-over point" exists where the theoretical speedup of quantum approaches is negated by the error 
rates inherent in current hardware constraints. By benchmarking across five distinct problem types (e.g., Max-Cut and Knapsack), we hypothesize 
that hybrid strategies—which combine the stability of classical logic with the state-space exploration of quantum mechanics—will emerge as the 
most viable path for bridging the current gap between quantum theory and real-world utility.

\section{Background}
To understand the comparative performance of classical and quantum optimization, one must first grasp the mechanics of Quantum Processing Units 
(QPUs) versus Central Processing Units (CPUs). While classical CPUs execute instructions linearly using bits (binary 0 or 1), quantum systems leverage 
qubits. Unlike bits, qubits utilize superposition, allowing them to represent a complex linear combination of both states simultaneously.

Two critical phenomena drive the theoretical power of the algorithms discussed in this work: entanglement and interference. Entanglement allows 
qubits to be perfectly correlated regardless of distance, while interference is used to amplify the probability of the "correct" solution while 
canceling out incorrect ones. However, current technology resides in the Noisy Intermediate-Scale Quantum (NISQ) era. In this stage, qubits are 
highly susceptible to decoherence—the loss of quantum information due to environmental interaction—and gate errors, which limit the circuit depth 
(the number of sequential operations) a computer can reliably perform.

This research specifically evaluates two primary quantum paradigms: Grover’s Algorithm and the Quantum Approximate Optimization Algorithm (QAOA). 
Grover’s provides a quadratic speedup for unstructured search by iteratively rotating the quantum state toward the target solution. Conversely, 
QAOA is a variational algorithm, a hybrid approach where a classical optimizer tunes quantum parameters to find the ground state of a Hamiltonian, 
typically representing a problem like Max-Cut or QUBO (Quadratic Unconstrained Binary Optimization). Understanding the trade-off between the 
theoretical "speedup" of these algorithms and the practical "noise floor" of NISQ hardware is essential for interpreting the benchmarking results 
provided in this paper.

\section{Related Works}
Even though quantum computing is relatively new, there has been a lot of work done in the field. The most relevant work to this project I found
was the paper by Airton Arruda, et al. titled "A Software-Based Comparative Study of Error Sensitivity in Classical and Quantum Optimization Algorithms"~\cite{aitor}.
This study compares three paradigms for solving the Weighted Max-Cut problem: a Genetic Algorithm, a Graph Neural Network, and a tensor 
network approach using Density Matrix Renormalization Group. Evaluating graphs from 10 to 250 nodes, it analyzes solution quality, runtime, 
and memory usage. Results show trade-offs in scalability, efficiency, and approximation performance. Additionally, papers such as 
"Introduction to Classical and Quantum Computing" by Thomas G. Wong~\cite{wong2019} and "Quantum Approximate Optimization Algorithm" by Leo Zhou, et al.~\cite{zhou2020} 
provide foundational insights into the principles and fundamentals to understand the algorithms used in this project. 

To understand the divergence between classical and quantum optimization, one must first establish the mechanical differences in how information 
is manipulated. Thomas G. Wong’s Introduction to Classical and Quantum Computing serves as a critical pedagogical anchor for this research. 
Wong’s work demystifies the transition from the deterministic logic of classical gates to the probabilistic nature of quantum operators. 
Specifically, his exposition on the Bloch Sphere and the mathematical representation of superposition provides the necessary context for why 
quantum algorithms can explore a search space more fluidly than their classical counterparts. By grounding the project in these fundamentals, 
we can better articulate why noise is not merely a "glitch" but a fundamental disruption of the delicate phase relationships that quantum speedups 
rely upon.Building upon these foundations, the work of Leo Zhou et al. in Quantum Approximate Optimization Algorithm (QAOA) introduces the specific 
mechanics of modern hybrid algorithms. QAOA is designed for the NISQ era, utilizing a variational approach where a classical optimizer tunes quantum 
parameters to minimize a cost function. Zhou’s analysis is vital because it highlights the theoretical potential of quantum circuits to find high-quality 
approximations for NP-hard problems, such as Max-Cut. However, Zhou also notes that the performance of QAOA is heavily dependent on the "depth" ($p$) of 
the circuit. This theoretical insight directly informs our investigation into the "coherence budget"—the point at which adding more quantum steps to 
improve a solution actually becomes counterproductive due to the accumulation of environmental noise.

While Wong and Zhou provide the theory, the work of Airton Arruda et al., A Software-Based Comparative Study of Error Sensitivity in Classical and 
Quantum Optimization Algorithms, offers a pragmatic blueprint for comparative methodology. Arruda’s study is particularly relevant because it moves 
beyond theoretical proofs to examine how algorithms behave when the "ideal" conditions of a blackboard are replaced by the limitations of a software-simulated environment.

Arruda et al. focus on the Weighted Max-Cut problem, comparing a Genetic Algorithm (GA), Graph Neural Networks (GNNs), and Density Matrix Renormalization Group (DMRG) 
approaches. Their findings regarding scalability are particularly instructive: while classical heuristics like GAs are robust and easy to scale to 250 nodes, they often 
settle into local optima. Conversely, tensor network approaches (DMRG) offer high precision but face significant memory bottlenecks as system complexity grows.

The most significant contribution of Arruda’s work to this project is the analysis of error sensitivity. By evaluating how solution quality degrades as noise models are 
introduced, Arruda identifies a "threshold of utility" for quantum-inspired methods. This research adopts and expands upon that logic, but pivots the focus toward Grover-based 
search and QAOA to see if these specific quantum structures exhibit the same sensitivity profiles as the tensor networks studied by Arruda. By comparing these against the "baseline" 
of a Genetic Algorithm—which effectively "ignores" quantum noise but struggles with search space dimensionality—we can pinpoint exactly where quantum methods lose their competitive edge.

Collectively, these works suggest that the future of optimization lies not in a total replacement of classical methods, but in a sophisticated understanding of their trade-offs. 
Wong provides the language, Zhou provides the tools, and Arruda provides the experimental rigour. This project builds on this lineage by creating a modular benchmarking framework 
that treats noise as a primary variable rather than an afterthought, ensuring that the comparison between a Greedy Search and a Quantum Algorithm is both fair and reflective of modern 
hardware realities.
\section{Methodology}
Dijkstra's algorithm was chosen for single-source shortest path computation on graphs with non-negative edge weights, following the implementation presented in CLRS~\cite{clrsAlgorithms}.

\section{Results}
An exemplary figure is shown in Figure~\ref{fig:dijkstra}. Any time a figure is used, it must have a capation and be directly referred to and discussed in the text.

%\begin{figure}[b] puts the image to the (closest) bottom of page
%\begin{figure}[t] puts the image to the (closest) top of page
%\begin{figure}[h] allows the image to float, putting it roughly wherever it was placed in the text.
\begin{figure}[b] 
\begin{center}
  \includegraphics[width=5cm]{imgs/dijkstra.png}
  \caption{Dijkstra's algorithm does not work on graphs with negative edge weights, as evidenced here.}
  \label{fig:dijkstra}
 \end{center}
\end{figure}



\section{Conclusion}

\bibliographystyle{acm}
\bibliography{bibliography.bib}

\end{document}
